\documentclass[leqno]{jsarticle}
\usepackage{amssymb}
\usepackage{amsthm}
\usepackage{amsmath}
\usepackage{comment}
	%可換図式用
		%\usepackage[all]{xy}
	%重くなるので使わいときはコメント化しておく。
%定理環境 thmはあまり使わずpropをなるべく使う。
%主張は基本的に証明の中で使い、そのため番号を付けない。
%注意環境を付けてみた。
\newtheorem{thm}{定理}
\newtheorem{prop}[thm]{命題}
\newtheorem{cor}[thm]{系}
\newtheorem{lem}[thm]{補題}
\newtheorem{df}[thm]{定義}
\newtheorem{exam}[thm]{例}
\newtheorem{fact}[thm]{事実}
\newtheorem*{claim}{主張}
\newtheorem*{cau}{注意}
\renewcommand{\proofname}{{\bf 証明}}
%矢印たち。
%\toは\rightarrowと同じ。\getsは\leftarrowと同じ。同値は\iff
%上向き矢はuparrow逆はdownarrow
\newcommand{\To}{\Longrightarrow}
\newcommand{\Gets}{\Longleftarrow}
%黒板太字
\newcommand{\nn}{\mathbb{N}}
\newcommand{\zz}{\mathbb{Z}}
\newcommand{\qq}{\mathbb{Q}}
\newcommand{\rr}{\mathbb{R}}
\newcommand{\cc}{\mathbb{C}}
%無限大
\newcommand{\mgn}{\infty}
%ギリシャン
\newcommand{\dpst}{\displaystyle}
\newcommand{\ep}{\varepsilon}
\newcommand{\al}{\alpha}
\newcommand{\be}{\beta}
\newcommand{\la}{\lambda}
\newcommand{\La}{\Lambda}
\newcommand{\ga}{\gamma}
%商集合
\newcommand{\qt}[2]{#1/\! #2}
%なんか演算
\newcommand{\card}{\text{  {\rm card}} }
\newcommand{\supp}{\text{  {\rm supp}} }
%なんか演算２
\newcommand{\bd}{\text{  {\rm Bdry}} }
\newcommand{\cl}{\text{  {\rm CL}} }
\newcommand{\nai}{\text{  {\rm Int}} }
\newcommand{\uu}{\mathcal{U}}
\newcommand{\vv}{\mathcal{V}}
\newcommand{\A}{\mathcal{A}}
\newcommand{\B}{\mathcal{B}}
\newcommand{\F}{\mathcal{F}}
%差集合は\setminusで統一する。等号の否定は\neq
%真部分集合は\subsetneqq　偏微分は\partial
%ディスプレイスタイル。\dfracでディスプレイ分数
%空集合は\Oを使う！！！！！！！！！！！！

\title{距離化可能定理：破片}
\author{電波通信}
\date{\today}
			\begin{document}
\maketitle

\section{準備}
\begin{thm}\label{pp}$T_3$空間に於いて
以下は同値
\begin{align}
&\mbox{$X$はパラコンパクト}\\ 
&\mbox{$X$の任意の開被覆は閉集合からなる局所有限な被覆によって細分される}\\
&\mbox{$X$の任意の開被覆は（閉とも開とも限らない集合からなる）局所有限な被覆によって細分される}
\end{align}
\end{thm}
\begin{proof}
$T_3$性から各点の閉近傍全体が基本近傍系を成してることから$(1)\To (2)$がわかる。また$(2)\To (3)$は明らかであるから以下の論理包含図式が得られる。
\[
(1)\To (2)\To (3)
\]
以下$(3)\To(1)$を証明する。\par
[$(3)\To (1)$の証明]\par
$\mathcal{V}$を$X$の開被覆とすると$(3)$から局所有限な（閉とも開とも限らない）被覆で細分されるそれを$\mathcal{A}$
と置く。$\mathcal{A}$は局所有限なので各$x$毎に$\mathcal{A}$の局所有限性を担保する$x$の近傍$W_x$が存在する。ここで$T_3$性の仮定からより強く$CL(W_x)$は$\mathcal{A}$の有限個の元としか交わらないとしてもよい。すると$\mathcal{W}=\{W_x\}_{x\in X}$は$X$の開被覆なので再び$(3)$から$\mathcal{W}$の局所有限な被覆である細分$\mathcal{B}$が存在する。この$\mathcal{B}$は$\mathcal{W}$の細分なので$\mathcal{B}$の任意の元の閉包$CL(B)$は$\mathcal{A}$の有限個の元としか交わらない。更に
\[
\mathcal{C}=\{CL(B)\ |\ B\in \mathcal{B}\}
\]
と置くとこれは局所有限な閉被覆で$C\in\mathcal{C}$は$\mathcal{A}$の有限個の元としか交わらない。
ここで$A\in \mathcal{A}$に対して
\[
U_A=X\setminus \bigcup\{C\ |\ C\cap A=\O\ C\in \mathcal{C}\}
\]
と定義すると$A\subset U_A$であり$\mathcal{C}$の局所有限性から$U_A$は開集合である。この時次が成り立つ。$C\in \mathcal{C}$について
\[
C\cap U_A\neq\O \iff C\cap A\neq\O
\]
なぜなら、$C\cap U_A\neq\O$とすると$U_A$の定義から$C\cap A\neq \O$でしかありえない。逆は明らか。
対偶を取れば
\[
C\cap U_A=\O \iff C\cap A=\O
\]
も成り立つ。よって$\mathcal{C}$は$\{U_A\}_{A\in \mathcal{A}}$の有限個の元としか交わらない。\par
$\{U_A\}_{A\in \mathcal{A}}$が局所有限であることを示そう。各点$x$について$\mathcal{C}$の局所有限性を担保する$x$の近傍$N$が存在する。$N$は$\mathcal{C}$の有限個の元
$C_1,C_2,\cdots ,C_m$としか交わらず、$\mathcal{C}$が被覆を成すことから
\[
x\in N\subset \bigcup_{i=1}^mC_i
\]
となる。よって$\{U_A\}_{A\in \mathcal{A}}$の元が$N$と交われば$C_1,C_2,\cdots ,C_m$のいづれかと必ず交わっているが、$C_1,C_2,\cdots ,C_m$の各々は$\{U_A\}_{A\in \mathcal{A}}$の有限個の元とし交わらないので$N$も$\{U_A\}_{A\in \mathcal{A}}$の有限個の元としか交わらない。よって$\{U_A\}_{A\in \mathcal{A}}$は局所有限である。最後に$\mathcal{A}$が$\mathcal{V}$の細分を成していることから$A\in \mathcal{A}$毎に
\[
A\subset O_A 
\]
となる$O_A\in \mathcal{U}$を選び、
\[
\{U_A\cap O_A\}_{A\in \mathcal{A}}
\]
とすると、明らかに$\mathcal{V}$の細分で$A\subset U_A\cap O_A$で$\mathcal{A}$が被覆を成してることから$\{U_A\cap O_A\}_{A\in \mathcal{A}}$も被覆で$\{U_A\}_{A\in \mathcal{A}}$が局所有限であることから$\{U_A\cap O_A\}_{A\in \mathcal{A}}$も局所有限である。以上で$X$がパラコンパクトであることが示された。
\par[$(3)\To (1)$の証明終わり]

\end{proof}

\begin{cor}\label{pa}
位相空間$X$に、各$A\in \mathcal{F}$が$X$の閉集合で部分空間としてパラコンパクトハウスドルフであるような局所有限な被覆$\mathcal{F}$が存在するならば$X$もパラコンパクトハウスドルフである。
\end{cor}
\begin{proof}
位相的な直和空間$Y=\dpst\coprod_{A\in \mathcal{F}}A$は各$A$が$T_4$\footnote{パラコンパクトハウスドルフ空が$T_4$になることを思い出せ。}なので$Y$も$T_4$になる。そして局所有限性から自然な写像
\[
i:\coprod_{A\in \mathcal{F}}A\to X
\]
は閉写像になる。$T_4$空間の閉像は$T_4$になるので$X$は$T_4$空間になり、特にハウスドルフ空間である。次にパラコンパクト性を示そう。\par
$\mathcal{U}$を$X$の開被覆とする。そして各$A\in \mathcal{F}$について
\[
\mathcal{U}_A=\{U\cap A\}_{U\in \mathcal{U}}
\]
と定義すると$\mathcal{U}_A$は$A$の開被覆である。よって$A$における$\mathcal{U}_A$の局所有限な細分である$A$の被覆が存在する。それを$\mathcal{R}_A$と置こう。そして
\[
\mathcal{R}=\{S\in \mathcal{R}_A\mid A\in \mathcal{F}\}=\bigcup_{A\in \mathcal{F}}\mathcal{R}_A
\]
と定義する。$\mathcal{F}$が被覆であることから$\mathcal{R}$は$X$の被覆となり、$\mathcal{U}$の細分となる。$\mathcal{R}$が局所有限であることを示そう。任意の$p\in X$について$\mathcal{F}$の局所有限性を担保する$p$の近傍$N$をとる。更に$\dpst O=\bigcap_{p\notin F,F\in \mathcal{F}}(X\setminus F)$とすると$\mathcal{F}$の局所有限性から$O$は$p$の属する開集合である。$V=N\cap O$は$p$の近傍である。そして$V$は$A_1,A_2,\dots A_s\in \mathcal{F}$以外の$\mathcal{F}$の元以外と交わらないとすると、$O$の定義から任意の$i=1,2,\dots,s$について$p\in A_i$である。各$\mathcal{U}_{A_i}$は$A_i$の中で局所有限なのでその局所有限性を担保する$p$の$A_i$における近傍$B_i$が存在する。$X$における$p$の近傍$U_i$が存在して$U_i\cap A_i=B_i$となる。さて、$\dpst W=V\cap \bigcap_{i=1}^sU_i$と置けばこれは$\mathcal{R}$の有限個の元としか交わらない。よって$\mathcal{R}$は局所有限である。そして先の定理\ref{pp}から$X$はパラコンパクトであることが分かる。($\mathcal{R}$の元は閉だとか開だとかであるとは一般に言えない。)
\end{proof}
先の系\ref{pa}の仮定における「局所有限」の部分を「有限」にした命題も成り立つ。系\ref{pa}の特別な場合に過ぎないが、一応別に述べておく。
\begin{cor}\label{pap}
位相空間$X$に、各$A\in \mathcal{F}$が$X$の閉集合で部分空間としてパラコンパクトハウスドルフであるような\textbf{有限な}被覆$\mathcal{F}$が存在するならば$X$もパラコンパクトハウスドルフである。
\end{cor}




\section{距離化可能：破片}
\begin{cau}
このセクション中でしばしば「距離空間の簡単な議論」と言うと思うが、それは$a$を中心とした半径$r$の開球を$B(a,r)$と書くときに
\[
p\in B(x,r)\To B(x,r)\subset B(p,2r)
\]
が成り立つと言うことである。
\end{cau}
\begin{fact}[Bing-長田-Smirnovの距離化可能定理]\label{dis}

$T_3$空間$X$について以下は同値
\begin{align}
&\mbox{$X$は距離化可能}\\
&\mbox{$X$は$\sigma$疎な開基を持つ}\\
&\mbox{$X$は$\sigma$局所有限な開基を持つ}
\end{align}
\end{fact}


\setcounter{equation}{0}
\begin{thm}\label{aaa}
ハウスドルフ空間に於いて以下は同値。
\begin{align}
&\mbox{パラコンパクトでかつ局所距離化可能}\\
&\mbox{距離化可能}
\end{align}
\end{thm}
\begin{proof}
$(2)\To(1)$は明らかである。逆を示そう。\par
まず$X$は$T_3$空間のである。そして
\[
\uu=\{V\ |\ \text{$V$は開集合で距離化可能}\}
\]
と定義すると$X$が局所距離化可能であることから$\uu$は被覆であり、$X$のパラコンパクト性から局所有限開被覆$\vv$で$\uu$を細分するものが存在する。特に$\vv$の任意の元は距離化可能である。$V$毎に$V$の位相と両立する距離を選び$d_V$としておく。さて各$n\in \nn$について
\[
\A_n=\{B(x,2^{-n};d_V)\ |\ x\in V\in \vv\}
\]
と定義し、$\B_n$を$\A_n$を細分する局所有限開被覆とする。ここで$x\in V$について
\[
B(x,r;d_V)=\{y\in V\ |\ d(x,y)<r\}
\]
であろ。さて
\[
\B=\bigcup_{n\in \nn}\B_n
\]
が$X$の基底になることを見よう。\par
$U$を$X$の任意の開集合とし$p\in U$とする。この時$p\in V$となる$V\in \vv$は有限個であり、そのような$V$について$U\cap V$は$V$の開集合であるので$n$を十分大きく取ると$p\in V$となる任意の$V\in \vv$について
\[
B(p,2^{-n};d_V)\subset V\cap U
\]
が成立する。そして$p\in G$となる$G\in \B_{n+1}$が存在し、さらに$\B_{n+1}$が$\A_{n+1}$の細分なので$G\subset B(y,2^{-n-1};d_W)$となる$y\in X$と$W\in \A_{n+1}$が存在する。$B(y,2^{-n-1};d_W)\subset W$なので$p\in W$であるから$B(p,2^{-n};d_W)\subset W\cap U$である。そして距離空間の簡単な議論から
\[
B(y,2^{-n-1};d_W)\subset B(p,2^{-n};d_W)\subset W\cap U
\]
なので結局
\[
p\in G\subset U
\]
が成り立つので$\B$は$X$の開基である。そして$\B$は$\sigma$局所有限なので$X$は事実\ref{dis}から距離化可能である。
\end{proof}


\setcounter{equation}{0}
\begin{lem}\label{bb}
位相空間$X$が有限個の距離化可能な閉部分空間の和なら$X$は距離化可能である。
\end{lem}
\begin{proof}
位相空間$X$が２つの距離化可能な閉部分空間の和であるときに証明すれば他の場合は帰納的に示される。\par
$X=A\cup B$で$A,B$は$X$の閉集合で距離化可能であるとする。$A,B$の位相と両立する距離をそれぞれ$d,e$とする。このとき$A,B$はパラコンパクトハウスドルフなので系\ref{pap}から$X$もパラコンパクトハウスドルフである。\par
$a\in A$について$P(a,r)=\{x\in A\ |\ d(a,x)<r\}$とし、$b\in B$について
$Q(b,r)=\{y\in B\ |\ e(b,y)<r\}$とする。そして$C(a,r),\ D(b,r)$をそれぞれ
\[
C(a,r)\cap A=P(a,r),\ D(a,r)\cap B=Q(b,r)
\]
となる$X$の開集合とし、$p\in \bd(A)\cap \bd(B)$について$O(p,r)=C(p,r)\cap D(p,r)$と置く。さらに$a\in \nai(A)$について$U(a,r)=C(a,r)\cap \nai (A)$とし、$b\in \nai (B)$について$V(b,r)=D(a,r)\cap \nai(B)$とする。このとき$O(x,r),\ U(a,r),\ V(b,r)$は$X$の開集合である。\par
以上を踏まえて
\[
\uu_n=\{O(p,2^{-n})\ |\ p\in \bd(A)\cap \bd(B)\}\cup \{U(a,2^{-n})\ |\ a\in \nai (A)\}\cup \{V(b,2^{-n})\ |\ b\in \nai(B)\}
\]
と定義し、$\vv_n$を$\uu_n$を細分する$X$の局所有限な開被覆とする。$X=\nai(A)\cup \nai (B)\cup (\bd(A)\cap \bd(B))$と直和分解されることに注意しよう。このとき
\[
\vv=\bigcup_{n\in \nn}\vv_n
\]
が$X$の開基になることを見よう。\par
$W$を$X$の任意の開集合とし、$x\in W$を任意に与える。\par
まず$x\in \bd(A)\cap\bd(B)$としよう。$W\cap A$は$A$の開集合で$W\cap B$は$B$の開集合なので十分大きく$n$を取れば$x\in P(x,2^{-n})\subset W,\ x\in Q(x,2^{-n})\subset W$となる。\par
さて$x\in G$となる$G\in \vv_{n+1}$を取れば$G\subset O(y,2^{-n-1})$となる$y\in \bd(A)\cap \bd(B)$が存在する。$O(y,2^{-n-1})$の定義から
\[
O(y,2^{-n-1})\subset P(y,2^{-n-1})\cup Q(y,2^{-n-1})
\]
となり、これと
距離空間の簡単な議論から$O(y,2^{-n-1})\subset P(x,2^{-n})\cup Q(x,2^{-n})\subset W$となる。よって$x\in G\subset W$となる$G\in \vv$が存在する。\par
次に$x\not\in \bd(A)\cap \bd(B)$のとき$N\in \nn$を十分大きく取れば$n\ge N$のとき、任意の\\
$p\in \bd(A)\cap \bd(B)$について$x\not\in O(p,2^{-n})$が成り立つように出来る。\par
さて、各$n\ge N$毎に$x\in G_n$となる$G_n\in \vv_n$を選んでおき、さらに
$G_n\subset E_n$となる$E_n\in \uu_n$も選んでおく。このとき$E_n$が$O(p,2^{-n})\ (p\in \bd(A)\cap \bd(B))$の形をしていることはあり得ないので$E_n$は$U(a,2^{-n})$もしくは$V(b,2^{-n})$の形をしている。\par
無限に多くの$\{E_n\}_{n\ge N}$の元が$U(a,2^{-n})$の形をしているときを考えよう。\par
このとき$N$をさらに十分大きくして$U(x,2^{-N})=P(x,2^{-N})\subset W$となるようにする。この時$N$より大きい$n$で$x\in G_n\subset E_n=U(a,2^{-n})=P(a,2^{-n})$となるものが存在する。すると距離空間の簡単な議論から
\[
x\in U(a,2^{-n})=P(a,2^{-n})\subset U(x,2^{-N})=P(x,2^{-N})\subset W
\]
となるので結局$x\in G_n\subset W$となる。無限に多くの$\{E_n\}_{n\ge N}$の元が$V(a,2^{-n})$の形をしているときも同様である。
\par
以上で$\vv$が開基となる事が分かった。そしてそして$\vv$は$\sigma$局所有限で$X$は$T_3$空間なので事実\ref{dis}から$X$は距離化可能である。
\end{proof}



\setcounter{equation}{0}
\begin{thm}
位相空間$X$に於いて以下は同値
\begin{align}
&\mbox{$X$は距離化可能な閉集合からなる局所有限な族の和で書ける。}\\
&\mbox{距離化可能}
\end{align}
\end{thm}
\begin{proof}
$(2)\To(1)$は明らかである。逆を示そう。\par
$\mathcal{F}$を距離化可能な閉集合からなる局所有限被覆としよう。$\mathcal{F}$の任意の元はパラコンパクトハウスドルフなので系\ref{pa}から$X$はパラコンパクトハウスドルフである。また任意の$p\in X$について$p$近の傍$N$が存在して$\mathcal{F}$の高々有限個の元としか交わらない。すると$N$は距離化可能な閉集合の有限個の和に含まれるので補題\ref{bb}から$N$は距離化可能である。つまり$X$は局所距離化可能である。よって定理\ref{aaa}から$X$は距離化可能であることが分かる。
\end{proof}
\begin{cau}
この補題は$A,B$に閉性を仮定しなければ成立しない。例えば非可算濃度の離散空間を１点コンパクト化した空間は無限遠点という１点と離散空間の和であり、それらは距離化可能であるが、１点コンパクト化された空間は明らかに距離化可能で無い。
\end{cau}


\setcounter{equation}{0}
\begin{thm}\label{asd}
正規空間$X$に於いて以下は同値。
\begin{align}
&\mbox{$X$は距離化可能な開集合からなる局所有限な族の和として書ける。}\\
&\mbox{距離化可能}
\end{align}
\end{thm}
\begin{proof}
$(2)\To(1)$は明らかである。逆を示そう。\par
$\uu$を距離化可能な開集合からなる局所有限な被覆としよう。すると$X$が正規であることから$\uu$の閉収縮$\{F_U\}_{U\in \uu}$が存在する。つまり、
\[
F_U\subset U
\]
となる閉集合の族$\{F_U\}_{U\in \uu}$が存在する。\par 
このとき$\{F_U\}_{U\in \uu}$は局所有限であり、各$F_U$は距離化可能であるから先の定理\ref{asd}から$X$は距離化可能であることが分かる。
\end{proof}

定理\ref{aaa}から以下の位相多様体の距離化可能性に関する定理は明らかであろう。
\begin{cor}
位相多様体$M$に於いて
\[
\mbox{距離化可能性$\iff$パラコンパクト性}
\]
\end{cor}










\begin{comment}
\section*{おまけ}
詳しく説明すると言ったが、次の命題を見れば分かるであろう。
\begin{prop}
$\F$を位相空間$X$の局所有限な閉被覆とすると包含写像から誘導される次の写像
\[
i:\coprod_{F\in \F}F\to X
\]
は閉写像である。ちなみにこの写像は１点の引き戻しが有限集合になるので完全写像でもある。
\end{prop}
\begin{proof}
$\dpst \coprod_{F\in\F}F$の閉集合$M$は$M_F\subset F$となる$X$の閉集合の族$\{M_F\}_{F\in \F}$と対応し、
\[
i(M)=\bigcup_{F\in \F}M_F
\]
となるが$\F$の局所有限性から$\{M_F\}_{F\in \F}$も局所有限になるから$\dpst \bigcup_{F\in \F}M_F$は閉集合であるから$i$は閉写像である。
\end{proof}
この事と、$T_1$及び、正規性、パラコンパクトハウスドルフ性が閉写像で保たれることを鑑みれば本文中の記述が良く分かることであろう。
\end{comment}




	
\begin{thebibliography}{9}
\bibitem{1}S.Willard,{\it General Topology},Dover Publications,2004,pp174,Problem 23G
\end{thebibliography}




			\end{document}



\begin{comment}
[可換図式]
\[\xymatrix{
&   & (2) \ar@{=>}[rd]&  &  &  & (5) \ar@{=>}[rd]&\\ 
& (1)\ar@{=>}[ru] \ar@{=>}[rd]&  &(4)\ar@{=>}[r]&(7)& (1)\ar@{=>}[ru] \ar@{=>}[rd]& & (7)&\\ 
&   & (3) \ar@{=>}[ur]&  &  &  &(6) \ar@{=>}[ur]&
}\]

[\bigin \endを使わない方法]

{\prop[aa] aaa}
で
\begin{prop}[aa]
aaa
\end{prop}
と同じ\df \thm \proof でも同様

[直和]
$\amalg$

[同値関係でよく使う波線]
\sim

[引数付きの新命令]
\newcommand{\prn}[1]{\|#1\|_p}

[番号のリセット]

\setcounter{equation}{0}


[字体の変更]

{\rm hogehoge}と使う。

\rm ローマン
\it イタリック
\sl 斜体

[supp]

\newcommand{\supp}{\text{{\rm supp}}}

[中央揃えの改行]

	\begin{eqnarray*}
		hoge1&=&hogehoge
		hoge2&=&hogehofe

	\end{eqnarray*}

&&の部分で揃えられる。






[場合分け]

	\begin{cases}
		hoge1 & \text{状況１}\\
		hoge2 & \text{状況２}\\
		・
		・
		・
	\end{cases}



\end{comment}





